\section{Proof of~\Cref{prop:cmulj-into-cmuljminus}}\label{app:proofs-of-sec-circular-mulj}
% 
% \subsection{Proof of~\Cref{prop:cmulj-into-cmuljminus}}
\label{subsec:proofs-of-cmulj-into-cmuljminus}



We give a bespoke combination of a negative translation and a Friedman-Dragalin `A translation' suitable for our purposes.
Writing $\N:= \mu X (1+X)$ and $\nega \N \sigma := \sigma \to \N$.
At the risk of ambiguity but to reduce syntax, we shall simply write $\neg \sigma$ for $\nega \N \sigma$ henceforth.\footnote{
Any confusion will be minimal as we shall not parametrise negation by any other A-translation in this work.} Also, we shall include the following derivable rules of $\cmuLJ$:
    \begin{equation}\label{eqn:negation-rules}
         \vlinf{\rr{\neg} }{}{\Gamma \seqar  \neg \sigma}{\Gamma, \sigma \seqar \N} 
 % & \dfn &  \vlinf{\rr{\to} }{}{\Gamma \seqar  \neg \sigma}{\Gamma \sigma \seqar \N}& 
 \qquad 
 \vlinf{\lr{\neg}}{}{\Gamma, \neg \sigma \seqar \N}{\Gamma \seqar \sigma}
%  & \dfn & \vlderivation{\vliin{\lr \to}{}{\Gamma, \neg \sigma \seqar \N}{\haxiom{\Gamma }{ \sigma}}{\axiom{\N}{\N}}}
% \\ \\ 
    \end{equation}
\[
     \vlinf{\rr{\neg \neg} }{}{\Gamma \seqar \neg \neg \sigma}{\Gamma \seqar \sigma}   
     % & \dfn  &
     % \vlderivation{
     % \vlin{\rr {\neg}}{}{\Gamma \seqar \neg \neg \sigma}
     %  {
     %  \vlin{\lr {\neg}}{}{\Gamma, \neg \sigma \seqar \N} 
     %   {
     %   \haxiom{\Gamma}{\sigma}
     %   }
     %  }
     % }
     % & \quad 
     \qquad 
     \vlinf{\lr {\neg \neg }}{}{\Gamma, \neg\neg \sigma \seqar \N}{\Gamma, \sigma \seqar \N }
     % & \dfn &
     % \vlderivation{
     % \vlin{\lr {\neg}}{}{\Gamma, \neg \neg \sigma \seqar \N}
     %   {
     %   \vlin{\rr  {\neg }}{}{\Gamma\seqar \neg \sigma}{\haxiom{\Gamma, \sigma}{\N}}
     %   }
     % }
\]


We define the translations $\negtrans \cdot$ and $\unnegtrans \cdot$ from arbitrary types to types over $\{\N,\times, \to , \mu\}$ as follows:
\[
% \begin{array}{r@{\ := \ }l}
%     \negtrans X & \neg \unnegtrans X\\
%     \negtrans 1 & \negn \unnegtrans 1 \quad \\
%     \negtrans{(\sigma \times \tau)} &\neg \unnegtrans{(  \sigma \times  \tau )}\\
%         \negtrans{(\sigma \to \tau)} &  \neg \unnegtrans{( \sigma \to  \tau)} \\
%     \negtrans{(\sigma + \tau)} & \neg \unnegtrans{(\sigma + \tau)}  \\
%      \negtrans{(\nu X \sigma)} &   \neg \unnegtrans{(\nu X. \sigma)}\\
%      \negtrans{(\mu X \sigma)} & \neg \unnegtrans{(\mu X. \sigma)}
% \end{array}
% \quad
\begin{array}{r@{\ := \ }l}
\negtrans \sigma & \neg \unnegtrans \sigma\\ 
\unnegtrans X & \negn X \\
    \unnegtrans 1 & \N \\
    \unnegtrans{(\sigma \times \tau)} & \negn (\negtrans \sigma \times \negtrans \tau) \\
    \unnegtrans{(\sigma \to \tau)} & \neg( \negtrans\sigma \to \negtrans \tau) \\
    \unnegtrans{(\sigma+\tau)} & \neg   \negtrans \sigma \times \neg \negtrans \tau \\
    \unnegtrans{(\nu X \sigma)} & \neg \neg \mu X \negn \negtrans\sigma[\negn X / X] \\
    \unnegtrans{(\mu X \sigma)} &  \neg \mu X \negtrans{\sigma}
\end{array}
\]

\begin{proposition}
    [Substitution]{\ }
    \todonew{what are the invariants? also mixed notations below, be consistent.}
    \begin{itemize}
    \item $\negtrans{(\sigma (\mu X. \sigma  ))}=\negtrans{\sigma}(\mu X. \negtrans{\sigma}) $
    \item $\negtrans{(\sigma (\nu X. \sigma))} = \negtrans{\sigma}(\neg \mu X. \neg \negtrans{\sigma}[\neg X/X])$
     \item $\unnegtrans{(\sigma (\nu X. \sigma))} = \unnegtrans{\sigma}(\neg \mu X. \neg \negtrans{\sigma}[\neg X/X])$.
        % \item  $\negtrans{\sigma(\tau)} = \negtrans \sigma (\negtrans \tau)$
        % \item $\unnegtrans {\sigma(\tau)}= \unnegtrans {\sigma} ({\neg \unnegtrans \tau} )$
        % \item $ \unnegtrans{ (\mu X. \sigma)} =  (\mu X. \negtrans \sigma) $
        % \item $\unnegtrans{(\nu X. \sigma)}= \neg \mu X.\neg {\negtrans \sigma(\neg X)}$
        % \item $\negtrans \sigma = \neg \unnegtrans \sigma$
    \end{itemize}
% .\oldtodo{clarify statement in obvious way}
\end{proposition}

We shall extend the notations $\negtrans \cdot$ and $\unnegtrans \cdot$ to cedents, e.g.\ writing $\negtrans {\Sigma}$ and $\unnegtrans {\Sigma}$, by distributing the translation over the list.
For a sequent $\Sigma \seqar \tau$, we also write $\negtrans{(\Sigma \seqar \tau)}$ for $\negtrans{\Sigma} , \unnegtrans \tau \seqar \N$





\begin{definition}
    [$N$-translation of steps]
    For each inference step \[\vliiinf{\rrule}{}{\Sigma \seqar \tau}{\Sigma_1 \seqar \tau_1}{\cdots }{\Sigma_n \seqar \tau_n}\] we define a gadget,
    \[
    \toks0={0.5}
    \vlderivation{
    \vltrf{\negtrans \rrule}{\negtrans{(\Sigma \seqar \tau)}}{\vlhy{\negtrans{(\Sigma_1 \seqar \tau_1)}}}{\vlhy{\cdots}}{\vlhy{\negtrans{(\Sigma_n \seqar \tau_n)}}}{\the\toks0}
    }
    \]
    as in~\Cref{fig:translation-inference-rules}.
    % \begin{itemize}
    %     \item \red{TODO: give definition}
    % \end{itemize}
       We lift this to a translation on coderivations $\der \mapsto \negtrans \der$ in the obvious way.
\end{definition}


\begin{figure*}
\centering
\[\small
    \begin{array}{rclrcl}
       \negtrans \ax  &\dfn & 
         \vlderivation{
         \vlin{\lr \neg }{}{\negtrans \sigma, \unnegtrans \sigma \seqar \N}{
         \axiom{\unnegtrans \sigma}{\unnegtrans \sigma}
           }
         }  
         &\qquad
         \negtrans \cut &\dfn &\vlderivation{
    \vliin{\cut}{}{\negtrans \Gamma, \negtrans \Delta, \unnegtrans \tau \seqar \N}
    {
    \vlin{\rr \neg }{}{\negtrans \Gamma \seqar \negtrans \sigma}{\haxiom{\negtrans \Gamma, \unnegtrans \sigma}{\N}}
    }
    {
    \haxiom{\negtrans \Delta, \negtrans \sigma, \unnegtrans \tau}{\N}
    }
    }
    \end{array}
    \]
    \bigskip
      \[\small
    \begin{array}{rclrcl}
    \negtrans {\rr \unit }  &  \dfn    & 
    \vlderivation{
    \vlin{\id}{}{\unnegtrans \unit \seqar \N}{\vlhy{}}
    }& \qquad 
\negtrans{\lr \unit}    &\dfn & 
\vlderivation{
\vlin{\wk}{}{\negtrans \Gamma, \negtrans \unit, \unnegtrans \sigma \seqar \N}
  {
\haxiom{\negtrans \Gamma, \unnegtrans \sigma}{\N}
  }
}
\end{array}
    \]
  \bigskip
      \[\small
    \begin{array}{rclrcl}
 \negtrans{\rr \to} & \dfn  &
    \vlderivation{
    \vlin{\lr \neg}{}{\negtrans \Gamma, \unnegtrans{(\sigma \to \tau)}\seqar \N}
     {
     \vlin{\rr \to}{}{\negtrans \Gamma \seqar \negtrans \sigma \to \negtrans \tau}
       {
       \vlin{\rr \neg}{}{\negtrans \Gamma, \negtrans \sigma \seqar \negtrans \tau}{\haxiom{\negtrans \Gamma, \negtrans \sigma, \unnegtrans \tau}{\N}}
       }
     }
    }
    &\qquad
\negtrans {\lr \to}    &\dfn & \vlderivation{
\vlin{\lr {\neg \neg}}{}{\negtrans \Gamma, \negtrans \Delta, \negtrans{(\sigma \to \tau)}, \unnegtrans \gamma \seqar \N}
    {
    \vliin{\lr \to}{}{\negtrans \Gamma, \negtrans \Delta, \negtrans \sigma \to \negtrans \tau , \unnegtrans \gamma \seqar \N}
   {
   \vlin{\rr \neg}{}{\negtrans \Gamma \seqar \negtrans \sigma}{\haxiom{\negtrans \Gamma, \unnegtrans \sigma}{\N}}
   }{
   \haxiom{\negtrans \Delta, \negtrans \tau, \unnegtrans \gamma}{\N}
   }
    }

}
\end{array}
    \]  
    
      \[\small
    \begin{array}{rclrcl}
\negtrans{\rr \times}&\dfn & \vlderivation{
\vlin{\lr \neg}{}{\negtrans \Gamma, \unnegtrans{(\sigma \times  \tau)}\seqar \N}
  {
  \vliin{\rr \times}{}{\negtrans \Gamma, \negtrans \Delta \seqar \negtrans \sigma \times \negtrans \tau}
    {
    \vlin{\rr \neg}{}{\negtrans \Gamma \seqar \negtrans \sigma}{\haxiom{\negtrans \Gamma, \unnegtrans \sigma}{\N}}
    }
    {
    \vlin{\rr \neg}{}{\negtrans \Delta \seqar \negtrans \tau}{\haxiom{\negtrans \Delta, \unnegtrans \tau}{\N}}
    }
  }
} &\qquad
\negtrans {\lr \times}&\dfn&\vlderivation{
\vlin{\lr {\neg \neg}}{}{\negtrans \Gamma, \negtrans{(\sigma \times \tau)}, \unnegtrans \gamma \seqar \N} 
    {
    \vlin{\lr \times }{}{\negtrans{\Gamma}, \negtrans \sigma \times \negtrans \tau , \unnegtrans \gamma \seqar \N} 
    {
    \haxiom{\negtrans \Gamma, \negtrans \sigma, \negtrans \tau, \unnegtrans \gamma}{\N}
    }
    }
  }
  \end{array}
    \]
      \bigskip
      \[\small
    \begin{array}{rclrcl}
\negtrans{\cev{\rr +}}&\dfn & \vlderivation{
\vlin{\lr \times }{}{\negtrans \Gamma, \unnegtrans{(\sigma + \tau)} \seqar \N}
  {
  \vlin{\wk}{}{\negtrans \Gamma, \neg  \negtrans \sigma  ,  \neg \negtrans \tau \seqar \N}
    {
    \vlin{\lr {\neg \neg}}{}{\negtrans \Gamma, \neg  \negtrans \sigma  \seqar \N}{\haxiom{\negtrans \Gamma, \unnegtrans \sigma}{\N}}
    }
  }
}
&
\qquad 
\negtrans{{+^1_r}}&\dfn & \vlderivation{
\vlin{\lr \times }{}{\negtrans \Gamma, \unnegtrans{(\sigma + \tau)} \seqar \N}
  {
  \vlin{\wk}{}{\negtrans \Gamma, \neg \negtrans \sigma  ,  \neg  \negtrans \tau \seqar \N}
    {
    \vlin{\lr {\neg \neg}}{}{\negtrans \Gamma, \neg  \negtrans \tau  \seqar \N}{\haxiom{\negtrans \Gamma, \unnegtrans \tau}{\N}}
    }
  }
}
\end{array}
\]
  \bigskip
\[
\negtrans {\lr +} \dfn\vlderivation{
\vlin{\lr {\neg  }}{}{\negtrans \Gamma, \negtrans{(\sigma + \tau)}, \unnegtrans \gamma \seqar \N} 
    {
    \vliq{\contr}{}{\negtrans{\Gamma}, \unnegtrans \gamma \seqar \neg  \negtrans \sigma \times \neg  \negtrans \tau}{
    \vliin{\lr \times }{}{\negtrans{\Gamma},\negtrans{\Gamma}, \unnegtrans \gamma \seqar \neg  \negtrans \sigma \times \neg  \negtrans \tau} 
    {
    \vlin{\lr \neg}{}{\negtrans \Gamma, \unnegtrans \gamma\seqar \neg \negtrans \sigma}{
    \haxiom{\negtrans \Gamma, \negtrans \sigma, \unnegtrans \gamma}{\N}
    }
    }{
    \vlin{\lr \neg}{}{\negtrans \Gamma, \unnegtrans \gamma\seqar \neg \negtrans \tau}{
    \haxiom{\negtrans \Gamma, \negtrans \tau, \unnegtrans \gamma}{\N}
    }
    }
    }
    }
}
    \]
      \bigskip
     \[\small
    \begin{array}{rclrcl}
    \negtrans {\rr \mu} & \dfn & 
    \vlderivation{
    \vlin{\lr {\neg }}{}{\negtrans \Gamma, \unnegtrans{(\mu X. \sigma)} \seqar \N}
    {
    \vlin{\rr \mu}{}{\negtrans \Gamma\seqar \mu X. \negtrans{\sigma}} 
    {
    \vlid{=}{}{\negtrans \Gamma\seqar  \negtrans{\sigma}(\mu X. \negtrans{\sigma})} 
      {
      \vlin{\rr {\neg }}{}{\negtrans \Gamma \seqar \negtrans{(\sigma (\mu X. \sigma))}}{
      \haxiom{\negtrans \Gamma, \unnegtrans{(\sigma (\mu X. \sigma))}}{\N}
      }
      }
    }
    }
    }
      & \qquad 
      \negtrans{\lr \mu}& \dfn & \vlderivation{
   \vlin{\lr {\neg \neg}}{}{\negtrans \Gamma, \negtrans{(\mu X. \sigma)} , \unnegtrans \gamma \seqar \N} 
       {
       \vlin{\rr \mu}{}{\negtrans \Gamma, \mu X. \negtrans{\sigma} , \unnegtrans \gamma \seqar \N}
         {
         \vlid{=}{}{\negtrans \Gamma,  \negtrans{\sigma}(\mu X. \negtrans{\sigma}) ,\unnegtrans \gamma \seqar \N}{
         \haxiom {\negtrans \Gamma,  \negtrans{(\sigma(\mu X.\sigma))} ,\unnegtrans \gamma}{  \N}
         }
         }
       }
   }
    \end{array}
    \]
       \bigskip
     \[\small
    \begin{array}{rclrcl}
    \negtrans{\rr \nu} & \dfn & 
    \vlderivation{
    \vlin{\lr {\neg \neg}}{}{\negtrans \Gamma, \unnegtrans{(\nu X. \sigma)} \seqar \N}{
    \vlin{\lr {\mu }}{}{\negtrans \Gamma,  \mu X. \neg {\negtrans \sigma(\neg X)}\seqar \N}
    {
    \vlin{\lr {\neg\neg }}{}{\negtrans \Gamma, \neg {\negtrans \sigma(\neg \mu X. \neg {\negtrans \sigma(\neg X)})} \seqar \N} 
    {
    \vlid{=}{}{\negtrans \Gamma,  {\unnegtrans \sigma(\neg \mu X. \neg {\negtrans \sigma(\neg X)})} \seqar \N} 
      {
      \haxiom{\negtrans \Gamma, \unnegtrans{(\sigma(\nu X. \sigma))}}{\N}
      }
    }
    }
    }
    }& \qquad
   \negtrans{\lr \nu}&\dfn &  \vlderivation{
   \vlin{\lr {\neg \neg }}{}{\negtrans \Gamma, \negtrans{(\nu X. \sigma )}, \unnegtrans \gamma \seqar \N} 
   {
   \vlin{\lr {\neg}}{}{\negtrans \Gamma,  \neg \mu X. \neg {\negtrans \sigma(\neg X)}, \unnegtrans \gamma \seqar \N} 
   {
   \vlin{\lr {\mu}}{}{\negtrans \Gamma, \unnegtrans \gamma \seqar  \mu X. \neg {\negtrans \sigma(\neg X)}} 
       {
       \vlin{\rr {\neg } }{}{\negtrans \Gamma, \unnegtrans \gamma \seqar  \neg {\negtrans \sigma(\neg \mu X. \neg {\negtrans \sigma(\neg X)})}}
         {
         \vlid{=}{}{\negtrans \Gamma,  \negtrans \sigma(\neg \mu X. \neg {\negtrans \sigma(\neg X)}) ,\unnegtrans \gamma \seqar \N}{
         \haxiom{\negtrans \Gamma,  \negtrans{(\sigma (\nu X. \sigma ))} ,\unnegtrans \gamma}{  \N}
         }
         }
       }
       }
       }
   }
    \end{array}
    \]
    \caption{Translation $\negtrans{(\_)}$ on inference rules.}
    \label{fig:translation-inference-rules}
\end{figure*}


\begin{proposition}\label{prop:translation-preserves-progressiveness}
        If $\der$ is progressing and regular then so is $\negtrans \der$.
\end{proposition}
\begin{proof}
    As for regularity, we observe that the translation maps any inference rule $\rrule$ to a gadget $\negtrans \rrule$ with constantly many rules. Concerning progressiveness, it is easy to check that:
    \begin{itemize}
        \item  for any infinite branch $(\Delta^j \seqar \tau^j)_j$ of $\der'$ there is an infinite branch $(\Sigma^i \seqar \sigma^i)_i$ of $\der$ such that $\Delta^{j_i} \seqar \tau^{j_i}= \negtrans{(\Sigma^i \seqar \sigma^i)}$ for some sequence  $j_0<j_1<\ldots<j_n<\ldots$
        \item  for any thread  in $\der$ connecting a formula $\gamma$ in $\Sigma^i \seqar \tau^i$ with a formula $\gamma'$ in $\Sigma^{i+1} \seqar \tau^{i+1}$ there is a thread  in $\der'$ connecting $\negtrans \gamma$ in $\negtrans{(\Sigma^i \seqar \tau^i)}$ with a formula $\negtrans{\gamma'}$ in $\negtrans{(\Sigma^{i+1} \seqar \tau^{i+1})}$; moreover, if the former thread has progressing points then the latter has. \qedhere
    \end{itemize}
\end{proof}


\begin{proposition}
\label{lem:computational-soundness}
If $\der_1 \cutrednorec \der_2$ then $\negtrans{\der}_1\cutrednorec^* \negtrans{\der}_2$ 
    % $\negtrans{\rrule}$ simulates $\rrule$ with respect to the following theory: \red{TODO: make this explicit here.}
\end{proposition}
\begin{proof}
The proof is routine. We just consider the cut-reduction step eliminating a cut $\rrule$ of the form $\rr\nu$ vs $\lr \nu$. W.l.o.g., we assume $\der$ ends with this cut. Then, $\der$ has the following shape:
\[
\vlderivation{
\vliin{\cut}{}{\Gamma, \Delta \seqar \tau}{\vlin{\rr \nu}{}{\Gamma \seqar \nu X. \sigma}{\vltr{\der_1}{\Gamma \seqar \sigma (\nu X. \sigma)}{\vlhy{\ \ }}{\vlhy{\ \ }}{\vlhy{\ \ }}}}{\vlin{\lr \nu}{}{\Delta, \nu X. \sigma  \seqar \tau}{\vltr{\der_2}{\Delta, \sigma (\nu X. \sigma)\seqar \tau}{\vlhy{\ \ }}{\vlhy{\ \ }}{\vlhy{\ \ }}}}
}
\]
Which is translated into a cut between the following two coderivations:
\[
\small
\vlderivation{
% \vliin{\cut}{}{\negtrans{\Gamma}, \negtrans \Delta \seqar \N }
% {
\vlin{\rr \neg}{}{\negtrans{\Gamma}\seqar \negtrans{(\nu X. \sigma)} }{
    \vlin{\lr {\neg \neg}}{}{\negtrans \Gamma, \unnegtrans{(\nu X. \sigma)} \seqar \N}{
    \vlin{\lr {\mu }}{}{\negtrans \Gamma,  \mu X. \neg {\negtrans \sigma(\neg X)}\seqar \N}
    {
    \vlin{\lr {\neg\neg }}{}{\negtrans \Gamma, \neg {\negtrans \sigma(\neg \mu X. \neg {\negtrans \sigma(\neg X)})} \seqar \N} 
    {
    \vlid{=}{}{\negtrans \Gamma,  {\unnegtrans \sigma(\neg \mu X. \neg {\negtrans \sigma(\neg X)})} \seqar \N} 
      {
      \vltr{\negtrans{\der_1}}{\negtrans \Gamma, \unnegtrans{(\sigma(\nu X. \sigma))}\seqar \N}{\vlhy{\ \ }}{\vlhy{\ \ }}{\vlhy{\ \ }}
      }
    }
    }
    }
    }
    }
    % }
     % {
     \]
     \[
     \vlderivation{
     \small
   \vlin{\lr {\neg \neg }}{}{\negtrans \Delta , \negtrans{(\nu X. \sigma )}, \unnegtrans \tau \seqar \N} 
   {
   \vlin{\lr {\neg}}{}{\negtrans \Delta,  \neg \mu X. \neg {\negtrans \sigma(\neg X)}, \unnegtrans \tau \seqar \N} 
   {
   \vlin{\lr {\mu}}{}{\negtrans \Delta, \unnegtrans \tau \seqar  \mu X. \neg {\negtrans \sigma(\neg X)}} 
       {
       \vlin{\rr {\neg } }{}{\negtrans \Delta, \unnegtrans \tau \seqar  \neg {\negtrans \sigma(\neg \mu X. \neg {\negtrans \sigma(\neg X)})}}
         {
         \vlid{=}{}{\negtrans \Delta,  \negtrans \sigma(\neg \mu X. \neg {\negtrans \sigma(\neg X)}) ,\unnegtrans \tau \seqar \N}{
         \vltr{\negtrans{\der_2}}{\negtrans \Delta,  \negtrans{(\sigma (\nu X. \sigma ))} ,\unnegtrans \tau\seqar  \N}{\vlhy{\ \ }}{\vlhy{\ \ }}{\vlhy{\ \ }}
         }
         }
       }
       }
       }
   % }
}
\]
  by applying a few cut-reduction  steps we obtain the following:
  \[
\vlderivation{
\vliin{\cut}{}{\negtrans \Gamma, \negtrans \Delta, \unnegtrans \tau \seqar \N}{
\vlin{\rr \neg}{}{\negtrans\Gamma\seqar  \negtrans{(\nu X. \sigma)} }{
\vltr{\negtrans{\der_1}}{\negtrans \Gamma, \unnegtrans{(\sigma (\nu X. \sigma))}\seqar \N}{\vlhy{\ \ }}{\vlhy{\ \ }}{\vlhy{\ \ }}
}
}
{
\vltr{\negtrans{\der_2}}{\negtrans \Delta , \negtrans{(\sigma (\nu X. \sigma))}, \unnegtrans \tau \seqar \N }{\vlhy{\ \ }}{\vlhy{\ \ }}{\vlhy{\ \ }}}
}
\]
Which is the translation of the following:
\[
\vlderivation{
\vliin{\cut}{}{\Gamma, \Delta \seqar \tau}
{\vltr{\der_1}{\Gamma \seqar \sigma (\nu X. \sigma)}{\vlhy{\ \ }}{\vlhy{\ \ }}{\vlhy{\ \ }}}
{\vltr{\der_2}{\Delta, \sigma (\nu X. \sigma)\seqar \tau}{\vlhy{\ \ }}{\vlhy{\ \ }}{\vlhy{\ \ }}}
} \qedhere
\]
\end{proof}


\begin{proposition}[(De)coding for $\N$]\label{prop:coding-and-decoding}
  There are  coderivations $\dercod: \N \seqar \negtrans \N$ and $\derdec: \negtrans \N \seqar \N$  over $\{\N,\to, \times, \mu\}$ in  $\cmuLJ$  such that, for any $n \in \Nat$:
  \[\small
\vlderivation{
\vliin{\cut}{}{\seqar \negtrans \N}{\vldr{\cod{n}}{\seqar \N}}{\vldr{\dercod}{\N \seqar \negtrans \N} }
}
 \cutrednorec^*
\vlderivation{\vldr{\negtrans{\cod n}}{\seqar \negtrans \N}}
\]
\[
\small
 \vlderivation{
\vliin{\cut}{}{\seqar \N}{\vldr{\negtrans{\cod n}}{\seqar \negtrans \N}}{\vldr{\derdec}{\negtrans \N \seqar  \N} }
}
 \cutrednorec^* 
\vlderivation{\vldr{\cod n}{\seqar \N}}
\]
\end{proposition}
\begin{proof}
$\dercod$ and $\derdec$ are defined, respectively, as follows:
    \[
% \vlderivation{ \vldr{\dercod}{\N \seqar \negtrans \N} } \qquad  \dfn \qquad 
\small
\vlderivation{
\vlin{\rr{\neg \neg}}{\bullet}{\N \seqar \negtrans{\N}}
  {
  \vlin{\mu r}{}{\N \seqar \mu X.(\neg  (\neg  \negtrans{\unit} \times  \neg \negtrans{ X}) )}
    {
    \vlin{\rr \neg}{}{\N \seqar \neg  (\neg  \negtrans{\unit} \times  \neg \negtrans{ \N}) ) }
      {
      \vlin{\lr \times}{}{\N ,  \neg  \negtrans{\unit} \times  \neg \negtrans{ \N} \seqar \N}{
      \vlin{\lr \mu}{}{\N ,  \neg  \negtrans{\unit} ,  \neg \negtrans{ \N} \seqar \N}
        {
        \vliin{\lr +}{}{ 1+     \N ,  \neg  \negtrans{\unit} ,  \neg \negtrans{ \N} \seqar \N}
          {
          \vliq{\wk}{}{ 1,  \neg  \negtrans{\unit} ,  \neg \negtrans{ \N} \seqar \N}
            {
            \vlin{\lr{\neg }}{}{   \neg \negtrans{1} \seqar \N}{\vlin{\rr \neg}{}{\seqar \negtrans 1}{\axiom{\N}{\N}}}
            }
          }
          {
          \vlin{\wk}{}{ \N ,  \neg  \negtrans{\unit} ,  \neg \negtrans{ \N} \seqar \N }
            {
            \vlin{\lr \neg}{}{\N, \neg \negtrans{\N}\seqar \N }{
            \vlin{\rr{\neg \neg}}{\bullet}{\N \seqar \negtrans \N}{\vlhy{\vdots}}
            }
            }
          }
        }
        }
      }
    }
  }
}
\]
\[
\small
% \vlderivation{\vldr{\derdec}{\negtrans \N \seqar \N}} \qquad  \dfn \qquad
\vlderivation{
\vlin{\lr {\neg \neg}}{\bullet}{\negtrans{\N} \seqar \N}
  {
  \vlin{\lr \mu}{}{ \mu X.(\neg  (\neg  \negtrans{\unit} \times  \neg \negtrans{ X}) )\seqar \N}
   {
   \vlin{\lr{\neg }}{}{\neg  (\neg  \negtrans{\unit} \times  \neg \neg \negtrans{\N}) \seqar \N}
      {
      \vliin{\rr \times}{}{\seqar \neg  \negtrans{\unit} \times  \neg \neg \negtrans{\N}}
        {
        \vlin{\rr {\neg}}{}{\seqar \neg \negtrans{1}}
         {
         \vlin{\wk}{}{\negtrans{1} \seqar \N}{\vlin{\rr N}{}{\seqar \N}{\vlhy{}}}
         }
        }{
        \vlin{\rr {\neg}}{}{\seqar \neg \negtrans{\N}}
         {
         \vliin{\cut}{}{ \negtrans{\N}\seqar \N}
         {
         \vlin{\lr{\neg \neg }}{\bullet}{\negtrans \N \seqar \N}{\vlhy{\vdots}}
         }
         {
         \vlin{\Nsucc}{}{\N \seqar \N}{\vlin{\id}{}{\N \seqar \N}{\vlhy{}}}
         }
         }
        }
      }
   }
  }
}
\]
where $0$ and $\succ$ are derivations as in~\Cref{fig:native-n-rules-mulj}.
\end{proof}

\begin{theorem}\label{thm:cmullminus-simulates-cmulj}
    If $\der: \vec \N \seqar \N$ in $\cmuLJ$ then there is  coderivation  $\der':\vec \N \seqar \N$ over $\{\N,\to, \times, \mu\}$ in $\cmuLJ$ such that $\der$ and $\der'$ represent the same functions on natural numbers.
\end{theorem}
\begin{proof}
    % \red{TODO: idea is to apply coding and decoding before and after $N$-simulation.}
Let $f$ be a numerical function of $\cmuLJ$, w.l.o.g.~we can assume $f$ is unary. This means that there is a derivation $\der$ of  $\N \seqar  \N$ such that 
\[\vlderivation{
\vliin{\cut}{}{\seqar \N}{\vldr{\cod{n}}{\seqar \N}}{\vldr{\der}{\N \seqar  \N} }
}
 \qquad  \cutrednorec^* \qquad 
\vlderivation{\vldr{\cod{f(n)}}{\seqar \N}}
\]
By Proposition~\ref{lem:computational-soundness} we have that, for any $n \in \Nat$:
 \[
 \begin{array}{rcl}
 \vlderivation{
\vliin{\cut}{}{\seqar \negtrans \N}{\vldr{\negtrans {\cod{n}}}{\seqar \negtrans \N}}{
\vlin{\rr \neg}{}{\negtrans\N \seqar   \negtrans\N}{
\vldr{ \negtrans \der}{ \negtrans\N, \unnegtrans \N \seqar   \N}
}
}
}
 &  \cutrednorec^* & 
\vlderivation{\vldr{ \negtrans{\cod{f(n)}}}{\seqar  \negtrans\N}}
\end{array}
\]
% \red{
% Since $\sem{\cut \ \dern{n} \ \der }= \sem{\cut} \ \sem{\dern n}\ \sem{\der}  \cutelime \sem{\der} \ \sem{\dern n}  \cutelime \sem{\dern {f(n)}}$, by applying~\Cref{lem:computational-soundness}, for all $\der': \seqar \semi{\N}$:
% \[
% \sem{\kol{(\cut \ \dern{n}\ \der)}}\sem{\der'} \cutelime  \sem{\kol{(\dern{f(n)})}}\sem{\der'}
% \]
% so  $\sem{\Kol{(\cut \ \dern{n}\ \der)}}\sem{\der'} \cutelime  \sem{\Kol{(\dern{f(n)})}}\sem{\der'}$ by~\Cref{eqn:big-k-small-k}. By \red{extensionality} this means that $\sem{\Kol{(\cut \ \dern{n}\ \der)}} \cutelime  \sem{\Kol{(\dern{f(n)})}}$,  but by definition:
% \[
% \sem{\Kol{(\cut \ \dern{n}\ \der)}}\cutelime  \sem{\cut \ \Kol{(\dern{f(n)})} \ \Kol{\der} } \cutelime \sem{\Kol \der} \sem{\Kol{(\dern n)}}
% \]
% }
Then, $f$ can be represented by the following coderivation $\der'$ over $\{\N,\to, \times, \mu\}$:
\[
\begin{array}{rcl}
    \vlderivation{
    \vliin{\cut}{}{\N \seqar \N}
    {
    \vliin{\cut}{}{\N \seqar \negtrans \N}{\vldr{\dercod}{\N \seqar \negtrans\N}}
    {
  \vlin{\rr \neg}{}{\negtrans\N \seqar   \negtrans\N}{
\vldr{ \negtrans \der}{ \negtrans\N, \unnegtrans \N \seqar   \N}
}
    }
    }
    {\vldr{\derdec}{\negtrans \N \seqar \N}}
    } \tag*{\qed}
\end{array}
\]
\renewcommand{\qed}{}
\end{proof}